\documentclass{beamer}
\usepackage{amsmath}
\usepackage{amssymb}

% Choose a theme
\usetheme{Madrid}
\usecolortheme{default}
\usefonttheme{serif}

% --- CUSTOM FOOTER ---
% Redefine the footline to remove author/institute and keep title/page number
\makeatletter
\setbeamertemplate{footline}
{
  \leavevmode%
  \hbox{%
  \begin{beamercolorbox}[wd=.5\paperwidth,ht=2.25ex,dp=1ex,center]{title in head/foot}%
    \usebeamerfont{title in head/foot}\insertshorttitle
  \end{beamercolorbox}%
  \begin{beamercolorbox}[wd=.5\paperwidth,ht=2.25ex,dp=1ex,right]{date in head/foot}%
    \usebeamerfont{date in head/foot}\insertframenumber{} / \inserttotalframenumber\hspace*{2ex}
  \end{beamercolorbox}}%
  \vskip0pt%
}
\makeatother
% --- END CUSTOM FOOTER ---

\title[Lecture 17: The spherical Earth]{Lecture 17: The spherical Earth}
\author{David Al-Attar}
\institute{Michaelmas term 2024}
\date{}

\begin{document}

% --- INTRODUCTION ---

\begin{frame}
    \titlepage
\end{frame}

\begin{frame}{Outline and Motivation}
    In this lecture, we apply ray theory to a spherically symmetric Earth model.
    \begin{itemize}
        \item We will explore how to locate an earthquake's \textbf{hypocentre} using seismic wave arrival times.
        \pause
        \item We will then derive the ray tracing equations for a spherical Earth, showing they can be greatly simplified.
        \pause
        \item Finally, we will discuss how these methods have been used to determine the Earth's internal velocity structure, and what the major features of that structure are.
    \end{itemize}
    \pause
    \begin{exampleblock}{Historical Context}
        The methods described are foundational to seismology. It's remarkable that much of our knowledge of Earth's deep structure was determined long before the advent of modern digital computers.
    \end{exampleblock}
\end{frame}

% --- HYPOCENTRE LOCATION ---

\begin{frame}{Locating an Earthquake's Hypocentre}
    \begin{block}{The Hypocentre}
        The point on a fault where an earthquake rupture initiates is known as the \textbf{hypocentre}. The first seismic waves recorded by distant seismometers originate from this point.
    \end{block}
    \pause
    \begin{itemize}
        \item We can use the arrival times of these first waves at a network of seismometers to determine the location ($ \mathbf{x}_h $) and origin time ($ t_h $) of the hypocentre.
        \pause
        \item For this, we need an accurate model of the Earth's seismic velocity structure.
        \pause
        \item \textbf{Ray theory} is used to efficiently calculate the travel time from a potential hypocentre to each seismometer.
    \end{itemize}
\end{frame}

\begin{frame}{The Inverse Problem for Location}
    The problem of finding the hypocentre is a classic inverse problem.
    \begin{itemize}
        \item The measured arrival time $ t_i $ at seismometer $i$ is related to the source parameters by:
        $$ t_{i}=t_{h}+T(\mathbf{x}_{i},\mathbf{x}_{h})+e_{i} $$
        where $T(\mathbf{x}_{i},\mathbf{x}_{h})$ is the ray-theoretic travel time and $e_i$ is measurement error.
    \end{itemize}
    \pause
    \begin{alertblock}{Least-Squares Minimisation}
        We have four unknowns ($t_h$ and the three components of $\mathbf{x}_h$). We estimate them by finding the values that minimise the sum of squared, weighted differences between observed and predicted arrival times:
        $$ J(\mathbf{x}_h, t_h) = \sum_{i=1}^{N_s} \frac{[t_h + T(\mathbf{x}_i, \mathbf{x}_h) - t_i]^2}{\sigma_i^2} $$
    \end{alertblock}
\end{frame}

\begin{frame}{Factors Affecting Accuracy}
    The accuracy of a hypocentre location depends on several factors:
    \begin{itemize}
        \item \textbf{Measurement Error:} Errors in picking arrival times on seismograms.
        \pause
        \item \textbf{Frequency Content:} The resolution is limited by the minimum wavelength. For seismic waves at 2 Hz, this is about 2.5 km.
        \pause
        \item \textbf{Station Geometry:} An even distribution of seismometers in distance and azimuth around the source gives the best constraint.
        \pause
        \item \textbf{Velocity Model Error:} Any inaccuracy in the assumed Earth structure will map directly into errors in the location.
    \end{itemize}
\end{frame}

% --- RAY THEORY IN SPHERICAL EARTH ---

\begin{frame}{Rays in a Spherical Earth: Setup}
    We now apply ray theory to an isotropic, spherically symmetric Earth, where velocity depends only on radius, $ v = v(r) $.
    \pause
    \begin{itemize}
        \item We start with Hamilton's equations, where the Hamiltonian for a P-wave is:
        $$ H(\mathbf{x}, \mathbf{p}) = \frac{1}{2} \alpha(r)^2 ||\mathbf{p}||^2 $$
        \item The ray path $\mathbf{x}(\sigma)$ and slowness $\mathbf{p}(\sigma)$ are given by:
        $$ \frac{d\mathbf{x}}{d\sigma} = \frac{\partial H}{\partial \mathbf{p}} \quad , \quad \frac{d\mathbf{p}}{d\sigma} = -\frac{\partial H}{\partial \mathbf{x}} $$
    \end{itemize}
    \pause
    \begin{alertblock}{Goal}
        Because of the spherical symmetry, this system of ODEs has a conserved quantity, which allows us to simplify the problem dramatically.
    \end{alertblock}
\end{frame}

\begin{frame}{A Conserved Quantity: The Ray Parameter}
    Due to the spherical symmetry, the quantity $ \mathbf{q} = \mathbf{x} \times \mathbf{p} $ is conserved along a ray.
    $$ \frac{d}{d\sigma}(\mathbf{x}\times \mathbf{p}) = \frac{d\mathbf{x}}{d\sigma}\times \mathbf{p} + \mathbf{x}\times\frac{d\mathbf{p}}{d\sigma} = 0 $$
    \pause
    \begin{itemize}
        \item This is analogous to the conservation of angular momentum.
        \item As a consequence, the ray path is always confined to the plane containing the Earth's center, the source, and the receiver.
    \end{itemize}
    \pause
    \begin{block}{The Ray Parameter}
        The magnitude of this conserved vector, $q$, is known as the \textbf{ray parameter}. It is constant for any given ray.
        $$ q = ||\mathbf{x} \times \mathbf{p}|| = r ||\mathbf{p}|| \sin \theta = \frac{r \sin\theta}{\alpha(r)} $$
        Here, $\theta$ is the angle between the position vector $\mathbf{x}$ and the ray direction $\mathbf{p}$.
    \end{block}
\end{frame}

\begin{frame}{Travel-Time Integrals}
    The conservation of the ray parameter $q$ allows us to find the travel time $T$ and epicentral distance $\Delta$ by direct integration, without needing to trace the full ray path.
    \pause
    \begin{itemize}
        \item A ray travelling down into the Earth will turn back towards the surface at a radius $r_t$ where it travels horizontally ($\theta=90^\circ$). At this turning point, $q = r_t / \alpha(r_t)$.
    \end{itemize}
    \pause
    \begin{alertblock}{Travel-Time and Distance Formulae}
        For a ray that turns at radius $r_t$ and returns to the surface at radius $a$:
        \begin{itemize}
            \item \textbf{Total Distance:} $ \Delta(q) = \int_{r_{t}}^{a}\frac{2\alpha(r)q}{r\sqrt{r^{2}-\alpha(r)^{2}q^{2}}}dr $
            \item \textbf{Total Time:} $ T(q) = \int_{r_{t}}^{a}\frac{2r}{\alpha(r)\sqrt{r^{2}-\alpha(r)^{2}q^{2}}}dr $
        \end{itemize}
    \end{alertblock}
\end{frame}

% --- EARTH STRUCTURE ---

\begin{frame}{Determining Earth's Velocity Structure}
    \begin{block}{The Inverse Problem}
        By measuring the travel times ($T$) of many seismic phases at various epicentral distances ($\Delta$) from thousands of earthquakes, seismologists have built up global travel-time curves. The inverse problem is to find the velocity structure $\alpha(r)$ that best explains these curves.
    \end{block}
    \pause
    \begin{exampleblock}{Key Historical Discoveries}
        This approach led to the discovery of Earth's major layers:
        \begin{itemize}
            \item \textbf{Richard Oldham} first identified the core.
            \item \textbf{Harold Jeffreys} argued for a fluid outer core.
            \item \textbf{Inge Lehmann} discovered the solid inner core.
        \end{itemize}
    \end{exampleblock}
\end{frame}

\begin{frame}{Overview of Earth's Velocity Structure (PREM)}
    The Preliminary Reference Earth Model (PREM) is a standard 1-D model of seismic velocities, density, etc.
    \begin{itemize}
        \item The general trend is that velocity increases with depth due to increasing pressure and density.
        \pause
        \item This smooth trend is interrupted by several sharp discontinuities.
    \end{itemize}
    \begin{center}
        % Placeholder for an image of the PREM model
        \textit{(A graph showing P-wave and S-wave velocity vs. depth would go here)}
    \end{center}
\end{frame}

\begin{frame}{Interpreting the Discontinuities}
    The major seismic discontinuities mark changes in composition or phase.
    \begin{itemize}
        \item \textbf{Core-Mantle Boundary (CMB):} A major compositional change from silicate rock (mantle) to liquid iron-nickel alloy (outer core). The fact that the outer core is liquid is confirmed by the inability of S-waves (which require shear strength) to travel through it.
        \pause
        \item \textbf{Inner-Core Boundary (ICB):} A phase change where the iron alloy freezes due to immense pressure, forming the solid inner core.
        \pause
        \item \textbf{Mantle Transition Zone (410 \& 670 km):} These are not compositional boundaries. They are solid-state phase transformations where mantle minerals rearrange into denser crystal structures under increasing pressure.
    \end{itemize}
\end{frame}

% --- SUMMARY ---

\begin{frame}{Summary: What You Need to Know}
    \begin{itemize}
        \item How an earthquake's \textbf{hypocentre} can be located by minimising the misfit between observed and predicted seismic arrival times.
        \pause
        \item In a spherically symmetric Earth, the \textbf{ray parameter} ($q$) is conserved along a ray, which simplifies the ray tracing problem to radial integrals.
        \pause
        \item How global travel-time curves are used as data in an inverse problem to determine the Earth's spherically averaged velocity structure.
        \pause
        \item The basic features of Earth's velocity structure: the crust, mantle, liquid outer core, and solid inner core, and the physical interpretation of the major discontinuities (CMB, ICB, 410, 670).
    \end{itemize}
\end{frame}

\end{document}
